% Ad Familiares 11.4 --- headnote
% ad-familiares-11-04, October 44 BC, Cisalpine Gaul

\textit{D. Junius Brutus Albinus to Cicero, from Cisalpine
Gaul between mid-October and the end of November 44 BC ---
Perseus dateline \textit{Scr. in Gallia citeriore inter
med. m. Oct. et ex. Nov. a. 710 (44)}. Decimus Brutus,
already \textit{imperator} and consul-designate for 42,
had occupied the province assigned him by Caesar and was
now campaigning against the Alpine tribes to season his
army and earn salutation as imperator in his own right. A
salutation conceded on the field by the troops gave him a
formal claim to a senatorial \textit{supplicatio} on his
return, and a stronger civil position against Antony, who
was already manoeuvring to drive him out of the province.}

\textit{The letter is a short, businesslike provincial
dispatch, and reads as the opening note of the active
correspondence between the two men. Decimus does not
quite ask for support: he takes Cicero's goodwill as
established (\textit{me tibi esse curae}), reports
victories, and closes with the practical request that
\textit{adiuva nos tua sententia} --- ``assist us with
your vote'' --- in the senatorial debate his own dispatch
to the Senate is about to provoke. The reply (11.5) and
its follow-up eleven days later (11.6) are Cicero's
acceptance of exactly that role.}
